%% 自動生成（tools/mknext.py）。各章の道しるべ「この章の先に」から、後ろの巻の章を集めたもの
\begin{longtblr}{colspec={Q[l,m,wd=58mm] X[l,m]}, width=\linewidth, rowhead=1, row{1}={bg=bkVol!14, font=\sffamily\bfseries\footnotesize},
  hlines={0.4pt,bkLine}, rowsep=1.5pt, cells={font=\footnotesize}}
この巻の章 & この道具を使う章 \\
{\bknextpill{7-2}\ 行列×ベクトル：形を変える機械} & {\bknextpill{8-3}\ 全微分：少し動いたときの変化} \\
{\bknextpill{7-5}\ 逆行列で連立方程式を解く} & {\bknextpill{9-10}\ 連立微分方程式を行列で：状態方程式} \\
{\bknextpill{7-7}\ 3次元の回転と座標の変換} & {\bknextpill{8-5}\ ベクトル場} \\
{\bknextpill{7-8}\ 固有値と固有ベクトル：向きの変わらない方向} & {\bknextpill{9-10}\ 連立微分方程式を行列で：状態方程式\\ \bknextpill{9-11}\ 行列の指数関数と固有値} \\
{\bknextpill{7-9}\ 対角化：まっすぐな方向に分けて考える} & {\bknextpill{9-10}\ 連立微分方程式を行列で：状態方程式\\ \bknextpill{9-11}\ 行列の指数関数と固有値} \\
{\bknextpill{7-10}\ 行列のくり返し：漸化式とオイラー法を行列で} & {\bknextpill{9-10}\ 連立微分方程式を行列で：状態方程式\\ \bknextpill{9-11}\ 行列の指数関数と固有値} \\
{\bknextpill{7-11}\ 行列の使いみち} & {\bknextpill{8-3}\ 全微分：少し動いたときの変化\\ \bknextpill{9-10}\ 連立微分方程式を行列で：状態方程式\\ \bknextpill{9-11}\ 行列の指数関数と固有値\\ \bknextpill{9-12}\ フーリエとラプラスの使いみち} \\
\end{longtblr}
